% % test.tex

% \begin{frame}
%     \begin{tikzpicture}

%         % % Cuadrícula de píxeles
%         % \draw[step=0.5cm,gray,very thin] (-2,-2) grid (2,2);

%         % % Píxeles resaltados
%         % \fill[red] (-1,0) rectangle (0,1);
%         % \fill[red!70!black] (0,0) rectangle (1,1);

%         % % Etiquetas de los colores
%         % \node at (-0.5,1.3) {\textbf{Color A}};
%         % \node at (0.5,1.3) {\textbf{Color B}};

%         % % Flecha de comparación
%         % \draw[->,thick] (0,0.5) -- (1,0.5);
%         % \node at (0.5,0.8) {$\Delta E$};

%         % Espacio de color LAB
%         \begin{axis}[
%                 width=10cm, height=10cm,
%                 axis lines=middle,
%                 % xlabel={$a^*$ (verde/rojo)}, ylabel={$b^*$ (azul/amarillo)}, zlabel={$L^*$ (luminosidad)},
%                 xmin=-10, xmax=60,
%                 ymin=-10, ymax=60,
%                 zmin=0, zmax=60,
%                 view={35}{30}
%             ]
%             % Puntos de los colores
%             \addplot3[only marks, color=red, mark=*] coordinates {(60, 50, 50)};
%             \addplot3[only marks, color=red!70!black, mark=*] coordinates {(40, 40, 30)};

%             % Línea de distancia entre ellos
%             \addplot3[dashed] coordinates {(60,50,50) (40,40,30)};

%             % Etiquetas
%             \node[anchor=south] at (axis cs:60,50,50) {A};
%             \node[anchor=north] at (axis cs:40,40,30) {B};

%         \end{axis}

%     \end{tikzpicture}
% \end{frame}

\begin{frame}

    \begin{tikzpicture}
        % % Plano de fondo (colores base)
        % \foreach \x in {0,1,2} {
        %         \foreach \y in {0,1,2} {
        %                 \fill[lightred] (\x,\y) rectangle (\x,\y) rectangle (\x+1,\y+1);
        %             }
        %     }

        % % Píxeles superiores (colores modificados)
        % \fill[colorTestA] (1,1) rectangle (2,2);
        % \fill[colorTestB] (2,1) rectangle (3,2);

        % % Flechas hacia los píxeles comparados
        % \draw[->, thick] (1.5,2) -- (1.5,2.5);
        % \draw[->, thick] (2.5,2) -- (2.5,2.5);

        % % Etiquetas de los colores
        % \node[above] at (1.5,2.5) {Color A};
        % \node[above] at (2.5,2.5) {Color B};
        \def\pointA{(40,50,60)}        % lab(60, 40, 50) rgb(225, 113, 56)
        \def\pointB{(25, 30,30)}        %  % lab(30, 25, 30) rgb(115, 53, 24)
        \def\pointC{(40,-50,40)}        % lab(40, 40, -50)rgb(112, 72, 177)
        % Espacio CIELAB
        \begin{axis}[
                % at={(5cm,0cm)},
                width=12cm, height=8cm,
                axis lines=middle,
                xlabel={$a^*$}, ylabel={$b^*$}, zlabel={$L^*$ (luminosidad)},
                xmin=-10, xmax=50,
                ymin=-50, ymax=50,
                zmin=0, zmax=60,
                tick label style={font=\tiny}, % Tamaño de las etiquetas
                label style={font=\tiny,gray}, % Tamaño de los nombres de los ejes
                % xticklabels={-50,-35,-20,...,50},
                % xtick={-50,-35,-20,...,50},
                yticklabels={-50,-35,-20,...,50},
                ytick={-50,-35,-20,...,50},
                zticklabels={0,20,...,60},
                ztick={0,20,...,60},
                % enlargelimits=5.1,
                % x label style={at={(axis description cs:70.5,-0.1)},anchor=north},
                view={45}{30}
            ]
            % Puntos de los colores
            \addplot3[only marks, color=colorTestA, mark=square*,mark options={scale=2.5}] coordinates {\pointA};
            \addplot3[only marks, color=colorTestB, mark=square*,mark options={scale=2.5}] coordinates {\pointB};
            \addplot3[only marks, color=colorTestC, mark=square*,mark options={scale=2.5}] coordinates {\pointC};
            \addplot3[only marks, color=gray, mark=*] coordinates {(40, 50, 0)}node[above right]{\scriptsize $A$};
            \addplot3[only marks, color=gray, mark=*] coordinates {(25, 30, 0)}node[above left]{\scriptsize $B$};
            \addplot3[only marks, color=gray, mark=*] coordinates {(40,-50,0)}node[above left]{\scriptsize $C$};

            % Línea de distancia entre ellos
            \addplot3[dashed] coordinates {\pointA \pointB}node[below right,colorA,pos=0.55]{\small $\Delta E_{ab}^*=30.7$};; % 30.7
            \addplot3[dashed] coordinates {\pointA \pointC}node[above left,colorA,pos=0.5] {\small $\Delta E_{ab}^*=54.8$};; % 54.8
            % Points dashed on xy plane
            \addplot3[thin,dashed,gray] coordinates {(0,50,0) (40,50,0)}; % A
            \addplot3[thin,dashed,gray] coordinates {(40,0,0) (40,50,0)}; % A
            \addplot3[thin,dashed,gray] coordinates {(40,50,0) (40,50,60)}; % A
            \addplot3[thin,dashed,gray] coordinates {(0,30,0) (25,30,0)}; % B
            \addplot3[thin,dashed,gray] coordinates {(25,0,0) (25,30,0)}; % B
            \addplot3[thin,dashed,gray] coordinates {(25,30,30) (25,30,0)}; % B
            \addplot3[thin,dashed,gray] coordinates {(0,-50,0) (40,-50,0)}; % C
            \addplot3[thin,dashed,gray] coordinates {(40,0,0) (40,-50,0)}; % C
            \addplot3[thin,dashed,gray] coordinates {(40,-50,0) (40,-50,40)}; % C
            % \def\pointA{(40,50,60)}        % lab(60, 40, 50) rgb(225, 113, 56)
            % \def\pointB{(25, 30,30)}        %  % lab(30, 25, 30) rgb(115, 53, 24)
            % \def\pointC{(40,-50,40)}        % lab(40, 40, -50)rgb(112, 72, 177)
            % \node[right,colorA] at (axis cs:45,35,50) {$\Delta E_{ab}^*=24.495$};

            % Etiquetas
            % \node[anchor=south,colorA] at (axis cs:30,50,60) {\small $A$}; % (50, 60, 40) rgb(217, 61, 54)
            % \node[anchor=north,colorA] at (axis cs:-35,-40,60) {\small $B$}; % (50,0,-40) rgb(54, 123, 187)
            % \node[anchor=north,colorA] at (axis cs:-50,40,60) {\small $C$};  % (50, -50, 40) rgb(21, 138, 44)
        \end{axis}
        % {$\Delta E_{ab}^*$}
    \end{tikzpicture}
\end{frame}

\thispagestyle{plain}
\begin{frame}

    \begin{tikzpicture}

        \begin{scope}[
                xshift=120,yshift=160, every node/.append style={yslant=0.5, xslant=-1},
                yslant=0.5, xslant=-1
            ]
            % Dibujar la malla de píxeles
            \foreach \x in {1} {
                    \foreach \y in {0,1,2} {
                            \fill[backgroundColor,fill opacity=0.35] (\x,\y) rectangle (\x+1,\y+1);
                        }
                }
            % Dibujar la cuadrícula con perspectiva
            \draw[step=1cm,gray,thin] (1,0) grid (2,3);
            \node[black] at(1.5,0.5){30};
            \node[black] at(1.5,1.5){25};
            \node[black] at(1.5,2.5){30};
        \end{scope}
        \begin{scope}[
                xshift=120,yshift=70, every node/.append style={yslant=0.5, xslant=-1},
                yslant=0.5, xslant=-1
            ]
            % Dibujar la malla de píxeles
            \foreach \x in {1} {
                    \foreach \y in {0,1,2} {
                            \fill[backgroundColor,fill opacity=0.35] (\x,\y) rectangle (\x+1,\y+1);
                        }
                }
            % Dibujar la cuadrícula con perspectiva
            \draw[step=1cm,gray,thin] (1,0) grid (2,3);
            \node[black] at(1.5,0.5){40};
            \node[black] at(1.5,1.5){40};
            \node[black] at(1.5,2.5){-50};
        \end{scope}
        % Definir el plano en perspectiva
        \begin{scope}[
                yshift=0, every node/.append style={yslant=0.5, xslant=-1},
                yslant=0.5, xslant=-1
            ]
            % Dibujar la malla de píxeles
            \foreach \x in {0,1,2,3,4} {
                    \foreach \y in {0,1,2,3,4} {
                            \fill[backgroundColor,fill opacity=0.5] (\x,\y) rectangle (\x+1,\y+1);
                        }
                }

            % Dibujar la silueta (ejemplo: una "L") 
            \foreach \x/\y in {1/1, 1/2, 1/3, 2/1, 3/1} {
                    \fill[colorTestB] (\x,\y) rectangle (\x+1,\y+1);
                }

            % Dibujar la cuadrícula con perspectiva
            \draw[step=1cm,gray,thin] (0,0) grid (5,5);
            \node[white] at (1.5,3.5){1};
            \node[white] at(1.5,2.5){2};
            \node[white] at(1.5,1.5){3};
            \node[white] at(2.5,1.5){4};
            \node[white] at(3.5,1.5){5};

        \end{scope}

        % Lineas entre planos
        % Dibujar líneas conectando puntos entre ambos planos
        % \draw[step=1cm,red,thin] (0,0) grid (5,5);
        % \draw[step=1cm,blue,thin] (-4,0) grid (0,5);
        \draw[colorA, thick,Triangle-] (-1,2.2) -- (-1,5);

        \begin{scope}[
                yshift=90, every node/.append style={yslant=0.5, xslant=-1},
                yslant=0.5, xslant=-1
            ]
            % Dibujar la malla de píxeles
            \foreach \x in {0,1,2,3,4} {
                    \foreach \y in {0,1,2,3,4} {
                            \fill[backgroundColor,fill opacity=0.5] (\x,\y) rectangle (\x+1,\y+1);
                        }
                }

            % Dibujar la silueta (ejemplo: una "L")
            \foreach \x/\y in {1/1, 1/2, 1/3, 2/1, 3/1} {
                    \fill[colorTestC] (\x,\y) rectangle (\x+1,\y+1);
                }

            % Dibujar la cuadrícula con perspectiva
            \draw[step=1cm,gray,thin] (0,0) grid (5,5);
            \node[white] at(1.5,3.5){1};
            \node[white] at(1.5,2.5){2};
            \node[white] at(1.5,1.5){3};
            \node[white] at(2.5,1.5){4};
            \node[white] at(3.5,1.5){5};

        \end{scope}

        \draw[Triangle-,thick,colorA](3.5,4.0)node[right]{%
            % $\begin{aligned}
            %         L^* & = 40 \\a^* & = 40 \\b^* & = -50\\
            %     \end{aligned}
            % $
        }
        to[out=180,in=90] (2.5,2.5);
        \draw[Triangle-,thick,colorA](3.6,7.2)node[right]{%
            % $\begin{aligned}
            %         L^* & = 30 \\a^* & = 25 \\b^* & = 30\\
            %     \end{aligned}
            % $
        }
        to[out=180,in=90] (2.5,5.5);

        % lab(30, 25, 30) rgb(115, 53, 24)
        % lab(40, 40, -50)rgb(112, 72, 177)

    \end{tikzpicture}
\end{frame}
\begin{frame}
    \frametitle{Diferencia de color \(\Delta E_{94}^*\)}



    La diferencia de color \(\Delta E_{94}^*\) entre dos colores en el espacio CIE Lab se calcula mediante la siguiente ecuación:

    \[
        \Delta E_{94}^* = \sqrt{
            \left( \frac{\Delta L^*}{k_L S_L} \right)^2 +
            \left( \frac{\Delta C^*}{k_C S_C} \right)^2 +
            \left( \frac{\Delta H^*}{k_H S_H} \right)^2
        }
    \]
    \scriptsize
    Donde:
    \begin{itemize}
        \item \(\Delta L^* = L_1^* - L_2^*\) es la diferencia en luminosidad.
        \item \(\Delta C^* = C_1^* - C_2^*\) es la diferencia en croma.
        \item \(\Delta H^*\) es la diferencia en tono, calculada como:
              \[
                  \Delta H^* = \sqrt{ (\Delta a^*)^2 + (\Delta b^*)^2 - (\Delta C^*)^2 }
              \]
        \item \(S_L\), \(S_C\) y \(S_H\) son factores de ponderación que dependen de la aplicación:
              \[
                  S_L = 1, \quad S_C = 1 + K_1 C_1^*, \quad S_H = 1 + K_2 C_1^*
              \]
        \item \(k_L\), \(k_C\) y \(k_H\) son constantes que dependen de las condiciones de observación:
              \[
                  k_L = k_C = k_H = 1 \quad \text{(para condiciones estándar)}
              \]
        \item \(K_1\) y \(K_2\) son constantes que dependen de la aplicación:
              \[
                  K_1 = 0.045, \quad K_2 = 0.015
              \]
    \end{itemize}


\end{frame}


% Título del póster
% \title{Recreation of the Mona Lisa using Genetic Algorithms and CIE Lab Color Metrics}
% \author{Your Name}
% \institute{Your University}
% \date{\today}



% Introducción
\begin{frame}{Introduction}
    \begin{itemize}
        \item \textbf{Objective}: Recreate the Mona Lisa using a genetic algorithm that optimizes color and structural similarity.
        \item \textbf{Motivation}: Explore the intersection of evolutionary computation and digital art.
        \item \textbf{Approach}: Use CIE Lab color metrics and genetic algorithms to generate images iteratively.
    \end{itemize}
\end{frame}

% Metodología
\begin{frame}{Methodology}
    \begin{columns}
        \begin{column}{0.5\textwidth}
            \textbf{Genetic Algorithm}:
            \begin{itemize}
                \item \textbf{Initialization}: Randomly generate a population of images.
                \item \textbf{Selection}: Choose the fittest individuals based on color similarity.
                \item \textbf{Crossover}: Combine features of two parents to create offspring.
                \item \textbf{Mutation}: Introduce small random changes to the offspring.
            \end{itemize}
        \end{column}
        \begin{column}{0.5\textwidth}
            \textbf{Color Metric}:
            \begin{itemize}
                \item Use \textbf{delta\_E\_CIE1994} to measure color difference:
                      \[
                          \Delta E_{94}^* = \sqrt{\left(\frac{\Delta L^*}{k_L S_L}\right)^2 + \left(\frac{\Delta C^*}{k_C S_C}\right)^2 + \left(\frac{\Delta H^*}{k_H S_H}\right)^2}
                      \]
                      where:
                      \begin{itemize}
                          \item \(\Delta L^*\), \(\Delta C^*\), \(\Delta H^*\): Differences in lightness, chroma, and hue.
                          \item \(k_L\), \(k_C\), \(k_H\): Weighting factors.
                          \item \(S_L\), \(S_C\), \(S_H\): Scaling factors.
                      \end{itemize}
            \end{itemize}
        \end{column}
    \end{columns}
\end{frame}

% % Esquema del algoritmo genético
% \begin{frame}{Genetic Algorithm Workflow}
%     \begin{center}
%         \begin{tikzpicture}[node distance=2cm]
%             \node (start) [startstop] {Start};
%             \node (init) [process, below of=start] {Initialize Population};
%             \node (eval) [process, below of=init] {Evaluate Fitness};
%             \node (select) [process, below of=eval] {Select Parents};
%             \node (crossover) [process, below of=select] {Crossover};
%             \node (mutate) [process, below of=crossover] {Mutate};
%             \node (newgen) [process, below of=mutate] {New Generation};
%             \node (stop) [decision, below of=newgen, yshift=-0.5cm] {Stop Condition?};
%             \node (end) [startstop, below of=stop, yshift=-0.5cm] {End};

%             \draw [arrow] (start) -- (init);
%             \draw [arrow] (init) -- (eval);
%             \draw [arrow] (eval) -- (select);
%             \draw [arrow] (select) -- (crossover);
%             \draw [arrow] (crossover) -- (mutate);
%             \draw [arrow] (mutate) -- (newgen);
%             \draw [arrow] (newgen) -- (stop);
%             \draw [arrow] (stop) -- node[anchor=east] {No} (eval);
%             \draw [arrow] (stop) -- node[anchor=south] {Yes} (end);
%         \end{tikzpicture}
%     \end{center}
% \end{frame}

% Resultados
\begin{frame}{Results}
    \begin{columns}
        \begin{column}{0.5\textwidth}
            \textbf{Generated Images}:
            \begin{itemize}
                \item Evolution of the Mona Lisa over generations.
                \item Comparison between the original and the final generated image.
            \end{itemize}
            % \includegraphics[width=\textwidth]{mona_lisa_evolution.png}
        \end{column}
        \begin{column}{0.5\textwidth}
            \textbf{Fitness Evolution}:
            \begin{itemize}
                \item Graph showing the improvement of fitness over generations.
            \end{itemize}
            % \includegraphics[width=\textwidth]{fitness_evolution.png}
        \end{column}
    \end{columns}
\end{frame}

% Discusión
\begin{frame}{Discussion}
    \begin{itemize}
        \item \textbf{Challenges}:
              \begin{itemize}
                  \item High computational cost.
                  \item Sensitivity to outliers in color metrics.
              \end{itemize}
        \item \textbf{Advantages}:
              \begin{itemize}
                  \item Effective convergence towards the target image.
                  \item Flexibility to use different color metrics.
              \end{itemize}
    \end{itemize}
\end{frame}

% Conclusiones
\begin{frame}{Conclusions}
    \begin{itemize}
        \item The genetic algorithm successfully recreates the Mona Lisa using CIE Lab color metrics.
        \item Future work includes:
              \begin{itemize}
                  \item Using more robust color metrics (e.g., CIE 2000).
                  \item Optimizing the algorithm for faster execution.
              \end{itemize}
    \end{itemize}
\end{frame}

% Referencias
\begin{frame}{References}
    \begin{itemize}
        \item \textbf{Colour Science}: \texttt{https://www.colour-science.org/}
        \item \textbf{Genetic Algorithms}: Goldberg, D. E. (1989). \textit{Genetic Algorithms in Search, Optimization, and Machine Learning}.
        \item \textbf{Image Processing}: Gonzalez, R. C., \& Woods, R. E. (2008). \textit{Digital Image Processing}.
    \end{itemize}
\end{frame}